\subsection{Plancherel's theorem}

We denote by A(G) the vector subspace of \(L^{1}(G)\) generated by the functions \(f * g\) for \(f, g \in L^{1}(G) \cap L^{2}(G)\).

PROPOSITION 7. --- The space A(G) is a self-adjoint ideal of L¹(G). It is contained in L¹(G) ∩ L²(G), and in the image of C(G) in L¹(G).

Let \(f \in \mathrm{L}^1(\mathrm{G})\). For every \(g \in \mathrm{L}^2(\mathrm{G})\), one has \(f * g \in \mathrm{L}^1(\mathrm{G})\). Therefore, the space \(\mathrm{L}^1(\mathrm{G}) \cap \mathrm{L}^2(\mathrm{G})\) is an ideal of \(\mathrm{L}^1(\mathrm{G})\), and the same is true of the space \(\mathrm{A}(\mathrm{G})\). The ideal \(\mathrm{A}(\mathrm{G})\) is self-adjoint.

Let \(f\) and \(g\) be in \(\mathrm{L}^2(\mathrm{G})\). The convolution product \(f * g\) is then the class of the continuous function given by

\[
y \mapsto \int_ {\mathrm{G}} f (y x ^ {- 1}) g (x) d x
\]

The second assertion follows from it.

Since \(\chi(f*g) = (\chi f)*( \chi g)\) for \(\chi \in \widehat{\mathrm{G}}\), \(f \in \mathrm{L}^1(\mathrm{G})\), and \(g \in \mathrm{L}^1(\mathrm{G})\) (INT, VIII, §3, no. 1, prop. 6), we have \(\chi h \in \mathrm{A}(\mathrm{G})\) for all \(h \in \mathrm{A}(\mathrm{G})\) and \(\chi \in \widehat{\mathrm{G}}\). Since \(\varepsilon_x * f = \gamma(x)f\) and the linear representation \(\gamma\) is isometric on \(\mathrm{L}^p(\mathrm{G})\) for all \(p\) (INT, VIII, §2, no. 5, prop. 8), we have \(\varepsilon_x * f \in \mathrm{A}(\mathrm{G})\) for all \(x \in \mathrm{G}\) and \(f \in \mathrm{A}(\mathrm{G})\).

We denote by \(\widehat{\mathrm{A}}(\mathrm{G})\) the image of \(\mathrm{A}(\mathrm{G})\) under the Fourier transform. It is a subspace of \(\mathcal{C}_0(\widehat{\mathrm{G}})\).

PROPOSITION 8. --- There exists a filter base \(\mathfrak{B}\) on \(\mathrm{A}(\mathrm{G}) \cap \mathcal{K}_{+}(\mathrm{G})\) such that the following conditions are satisfied:

\begin{enumerate}
\def\labelenumi{(\roman{enumi})}
\item
  For every element \(\varphi\) of a member of \(\mathfrak{B}\), one has \(\|\varphi\|_1 = 1\) and \(\|\mathcal{F}(\varphi)\|_{\infty} \leqslant 1\);
\item
  We have
\end{enumerate}

\[
\lim _ {\varphi , \mathfrak {B}} \varphi \cdot d x = \varepsilon_ {e}
\]

in the space \(\mathcal{C}'(\mathrm{G})\) of measures with compact support on G, endowed with the topology of uniform convergence on compact subsets of \(\mathcal{C}(\mathrm{G})\);

\begin{enumerate}
\def\labelenumi{(\roman{enumi})}
\setcounter{enumi}{2}
\tightlist
\item
  We have
\end{enumerate}

\[
\lim _ {\varphi , \mathfrak {B}} \mathcal {F} (\varphi) = 1
\]

for the topology of compact convergence on \(\widehat{\mathbf{G}}\) ;

\begin{enumerate}
\def\labelenumi{(\roman{enumi})}
\setcounter{enumi}{3}
\tightlist
\item
  For \(p = 1\) or \(p = 2\), and for every \(f \in \mathrm{L}^p(\mathrm{G})\), one has \(\varphi * f \in \mathrm{A}(\mathrm{G})\) for every \(\varphi\) belonging to a member of \(\mathfrak{B}\), and
\end{enumerate}

\[
\lim _ {\varphi , \mathfrak {B}} \varphi * f = f
\]

in L \(^{p}\) (G).

Let \(K_{0}\) be a fixed compact neighborhood of e in G. Let \(B_{0}\) be a base of the neighborhood filter of e in G consisting of compact symmetric neighborhoods contained in \(K_{0}\) (cf.~TG, III, p.~4). For \(K \in B_{0}\), let \(X_{K}^{\prime}\) be the set of functions \(\psi \in \mathcal{K}_{+}(G)\) such that \(\operatorname{Supp}(\psi) \subset K\) and \(\int \psi(x)dx = 1\); it is nonempty (Lemma 1 of II, p.~200). Let \(X_{K}\) be the set of functions \(\psi * \psi\) for \(\psi \in X_{K}^{\prime}\). It is nonempty and contained in \(A(G) \cap \mathcal{K}_{+}(G)\). The set \(\mathfrak{B}\) whose elements are the sets \(X_{K}\) for K varying in \(B_{0}\) is a filter base on \(A(G) \cap \mathcal{K}_{+}(G)\). We prove that \(\mathfrak{B}\) satisfies the required properties.

If \(X \in \mathfrak{B}\) and \(\varphi \in X\), one has \(\| \varphi \|_1 = \int_{G} \varphi(x) dx = 1\), hence \(\| \mathcal{F}(\varphi) \|_\infty \leqslant 1\), which establishes property (i).

Property (ii) follows from INT, VIII, § 2, no. 7, corollary 1 of lemma 4. A compact subset of \(\widehat{\mathbf{G}}\) is a compact subset of \(\mathcal{C}(\mathbf{G})\), hence (ii) implies \(\lim_{\varphi, \mathfrak{B}} \mathcal{F}(\varphi) = 1\) for the topology of compact convergence on \(\widehat{\mathbf{G}}\), that is, (iii).

Finally, let \(p=1\) or \(p=2\). Let \(f\in\mathrm{L}^{p}(\mathrm{G})\). We have \(\varphi*f\to f\) in \(\mathrm{L}^{p}(\mathrm{G})\) according to the filter \(\mathfrak{B}\) (INT, VIII, §4, no. 7, prop. 20). Moreover, for every K in \(B_{0}\) and \(\varphi\in X_{K}\), there exists \(\psi\in X_{K}^{\prime}\) such that \(\varphi=\psi*\psi\), whence \(\varphi*f=\psi*(\psi*f)\). We have \(\psi\in L^{1}(G)\cap L^{2}(G)\) and \(\psi*f\in L^{1}(G)\cap L^{2}(G)\), hence \(\varphi*f\in A(G)\).

COROLLARY 1. --- The space A(G) is dense in L \(^{1}\) (G) and in L \(^{2}\) (G). It is also dense in Stell(G), and its image \(\widehat{\mathrm{A}}(G)\) under the Fourier transformation is dense in \(\mathcal{C}_{0}(\widehat{\mathrm{G}})\).

Assertion (iv) of the proposition provides the first assertion. Since L \(^{1}\) (G) is dense in Stell(G), the second follows from it, and the last then follows from prop. 5 of II, p.~209.

COROLLARY 2. --- For \(f \in \mathrm{A}(\mathrm{G})\), let \(\Omega_f\) be the set of \(\chi \in \widehat{\mathrm{G}}\) such that \(\mathcal{F}(f)(\chi) \neq 0\). The sets \(\Omega_f\) form an open covering of \(\widehat{\mathrm{G}}\).

This follows from the preceding corollary since, for every \(\chi \in \widehat{\mathrm{G}}\), the map \(f \mapsto \mathcal{F}(f)(\chi)\) is a nonzero character of \(\mathrm{L}^1(\mathrm{G})\).

Recall that the left regular representation \(\pmb{\gamma}\) of Stell(G) on \(\mathrm{L}^2 (\mathrm{G})\) (cf.~I, p.~125, no. 13) is denoted \(\pmb{\gamma}(\varphi)f = \varphi *f\) for \(\varphi \in \mathrm{Stell}(\mathrm{G})\) and \(f\in \mathrm{L}^2 (\mathrm{G})\).

Lemma 3. --- For every \(f \in \mathrm{A}(\mathrm{G})\), there exists a unique bounded measure \(\mu_{f}\) on \(\widehat{\mathrm{G}}\) such that

\[
(\varphi * f) (e) = \int_ {\widehat {G}} \mathcal {F} (\varphi) d \mu_ {f}\tag{17}
\]

for all \(\varphi\in\mathrm{Stell}(\mathrm{G})\).

Moreover, for all \(f\) and \(g\) in A(G), we have the equality

\[
\mathcal {F} (f) \cdot \mu_ {g} = \mathcal {F} (g) \cdot \mu_ {f},\tag{18}
\]

between the measure with density \(\mathcal{F}(f)\) with respect to \(\mu_g\) and the measure with density \(\mathcal{F}(g)\) with respect to \(\mu_f\).

Let \(f\) and \(g\) be elements of \(\mathrm{L}^{1}(\mathrm{G}) \cap \mathrm{L}^{2}(\mathrm{G})\). For every \(\varphi \in \mathrm{Stell}(\mathrm{G})\), one has \(\varphi * f \in \mathrm{L}^{2}(\mathrm{G})\) and \(\| \varphi * f \|_{2} \leqslant \| \varphi \|_{*} \| f \|_{2}\) (I, p.~126, formula (8)). Moreover, we have \(\varphi * (f * g) = (\varphi * f) * g\) (loc. cit., formula (9)). This latter function belongs to the closure \(\mathcal{C}_{0}(\mathrm{G})\) of \(\mathcal{K}(\mathrm{G})\) in \(\mathcal{C}(\mathrm{G})\) (INT, VIII, §4, n° 5, prop. 15). Moreover, we have

\[
\left\| \varphi * (f * g) \right\| _ {\infty} \leqslant \left\| \varphi * f \right\| _ {2} \| g \| _ {2} \leqslant \left\| \varphi \right\| _ {*} \| f \| _ {2} \| g \| _ {2}.
\]

Since the functions \(f * g\) for \(f\) and \(g\) in \(\mathrm{L}^1(\mathrm{G}) \cap \mathrm{L}^2(\mathrm{G})\) generate \(\mathrm{A}(\mathrm{G})\), one deduces that \(\varphi * f \in \mathcal{C}_0(\mathrm{G})\) for all \(f \in \mathrm{A}(\mathrm{G})\) and \(\varphi \in \mathrm{Stell}(\mathrm{G})\), and that the map \(\varphi \mapsto (\varphi * f)(e)\) is a continuous linear form on Stell(G). Since \(\mathcal{F}\) is an isomorphism from Stell(G) onto \(\mathcal{C}_{0}(\widehat{\mathrm{G}})\) (prop. 5 of II, p.~209), the first assertion follows from it.

Now let \(f\) and \(g\) be in \(\mathrm{A}(\mathrm{G})\). For \(\varphi \in \mathrm{L}^1 (\mathrm{G})\), one has

\[
\begin{array}{r l r} & & {\big (\mathcal {F} (f) \cdot \mu_ {g} \big) (\mathcal {F} (\varphi)) = \int_ {\widehat {\mathrm{G}}} \mathcal {F} (\varphi) \mathcal {F} (f) d \mu_ {g} = \int_ {\widehat {\mathrm{G}}} \mathcal {F} (\varphi * f) d \mu_ {g}} \\ & & {= ((\varphi * f) * g) (e).} \end{array}\tag{19}
\]

Since \((\varphi*f)*g=(\varphi*g)*f\) and since the image of \(L^{1}(G)\) under the Fourier transform is dense in \(\mathcal{C}_{0}(\widehat{\mathrm{G}})\) (cor. of II, p.~210), one deduces from formula (19) that formula (18) is satisfied for all \(f\) and \(g\) in \(A(G)\).

Lemma 4. --- There exists a unique measure \(\nu\) on \(\widehat{\mathbf{G}}\) such that

\[
\mu_ {f} = \mathcal {F} (f) \cdot \nu
\]

for all \(f \in \mathrm{A}(\mathrm{G})\). For \(f \in \mathrm{A}(\mathrm{G})\), one has \(\mathcal{F}(f) \in \mathrm{L}^{1}(\widehat{\mathrm{G}}, \nu) \cap \mathrm{L}^{2}(\widehat{\mathrm{G}}, \nu)\). Let \(f \in \mathrm{A}(\mathrm{G})\). Denote by \(\Omega_{f}\) the open set in \(\widehat{G}\) consisting of the \(\chi \in \widehat{G}\) such that \(\mathcal{F}(f)(\chi) \neq 0\). Let \(\varphi\) be the characteristic function of \(\widehat{G} - \Omega_{f}\). For every \(g \in \mathrm{A}(\mathrm{G})\), we then have

\[
\int_ {\widehat {\mathrm{G}}} \mathcal {F} (g) d (\boldsymbol {\varphi} \cdot \boldsymbol {\mu} _ {f}) = \int_ {\widehat {\mathrm{G}}} \varphi \mathcal {F} (f) d \mu_ {g} = 0
\]

given formula (18).

By Corollary 1 of II, p.~212, the image \(\widehat{\mathrm{A}}(\mathrm{G})\) of \(\mathrm{A}(\mathrm{G})\) under the Fourier transform is dense in \(\mathcal{C}_0(\widehat{\mathrm{G}})\). We then deduce from the preceding formula that \(\varphi \cdot \mu_f = 0\), so that \(\mu_f\) is concentrated on \(\Omega_f\) (INT, IV, §4, no. 7, def. 4). Let \(\nu_f\) be the measure on \(\Omega_f\) with density \(\mathcal{F}(f)^{-1}\) with respect to \(\mu_f|\Omega_f\).

The sets \(\Omega_f\), for \(f \in \mathrm{A}(\mathrm{G})\), form an open covering of \(\widehat{\mathrm{G}}\) (cor. 2). For all \(f\) and \(g\) in \(\mathrm{A}(\mathrm{G})\), formula (18) shows that \(\nu_f|(\Omega_f \cap \Omega_g) = \nu_g|(\Omega_f \cap \Omega_g)\). Consequently, there exists a unique measure \(\nu\) on \(\widehat{\mathrm{G}}\) such that \(\nu_f = \nu|\Omega_f\) for all \(f \in \mathrm{A}(\mathrm{G})\) (INT, III, §2, no. 1, prop. 1).

If \(f \in \mathrm{A}(\mathrm{G})\), the measures \(\mu_f\) and \(\mathcal{F}(f) \cdot \nu\) are concentrated on \(\Omega_f\), and their restrictions to \(\Omega_f\) are equal to \(\mathcal{F}(f) \cdot \nu_f\); these measures are therefore equal.

Since \(\mu_f\) is a bounded measure, the Fourier transform \(\mathcal{F}(f)\) belongs to the space \(\mathrm{L}^1 (\widehat{\mathrm{G}},\nu)\). Moreover, since \(\mathcal{F}(f)\) belongs to \(\mathcal{C}_0(\widehat{\mathrm{G}})\), we also have \(\mathcal{F}(f)\in \mathrm{L}^2 (\widehat{\mathrm{G}},\nu)\).

Formula (17) now reads, for \(\varphi\in\mathrm{Stell}(\mathrm{G})\) and \(f\in\mathrm{A}(\mathrm{G})\):

\[
(\varphi * f) (e) = \int_ {\widehat {G}} \mathcal {F} (\varphi) \mathcal {F} (f) d \nu .\tag{20}
\]

In particular, for \(f\) and \(g\) in \(\mathrm{A}(\mathrm{G})\), one has

\[
\int_ {\widehat {\mathrm{G}}} \mathcal {F} (f) \mathcal {F} (g) d \nu = (f * g) (e) = \int_ {\mathrm{G}} f (x) g (x ^ {- 1}) d x.\tag{21}
\]

PROPOSITION 9. — The measure \(\nu\) characterized by Lemma 4 is a Haar measure on \(\widehat{\mathbf{G}}\).

Let \(\chi \in \widehat{\mathrm{G}}\). For \(f\) and \(g\) in \(\mathrm{A}(\mathrm{G})\), let us apply formula (21) to \(\chi f\) and \(\chi g\). The right-hand side is unchanged, hence

\[
\nu (\mathcal {F} (f) \mathcal {F} (g)) = \nu (\mathcal {F} (\chi f) \mathcal {F} (\chi g)).
\]

From formulas (12) of II, p.~208 and (13) of II, p.~208, it follows that

\[
\nu (\mathcal {F} (f) \mathcal {F} (g)) = \nu (\varepsilon_ {\chi} * (\mathcal {F} (f) \mathcal {F} (g))).
\]

We then deduce from INT, VIII, §4, n° 3, prop. 7 that

\[
\nu (\mathcal {F} (f) \mathcal {F} (g)) = (\varepsilon_ {\chi^ {- 1}} * \nu) (\mathcal {F} (f) \mathcal {F} (g)),
\]

that is,

\[
(\mathscr {F} (f) \cdot \nu) (\mathscr {F} (g)) = (\mathscr {F} (f) \cdot (\varepsilon_ {\chi^ {- 1}} * \nu)) (\mathscr {F} (g)).
\]

As the space \(\widehat{\mathrm{A}}(\mathrm{G})\) is dense in \(\mathcal{C}_{0}(\widehat{\mathrm{G}})\) (cor. 1), it follows that equality holds

\[
\mathscr {F} (f) \cdot \nu = \mathscr {F} (f) \cdot (\varepsilon_ {\chi^ {- 1}} * \nu).
\]

The measures \(\nu\) and \(\varepsilon_{\chi^{-1}} * \nu\) therefore coincide on the open set \(\Omega_f\) where \(\mathcal{F}(f)\) is nonzero. By Corollary 2 and INT, III, §2, no. 1, cor. of Prop. 1, these measures are therefore equal.

This shows that the measure \(\nu\) is proportional to a Haar measure on \(\widehat{\mathrm{G}}\). Let \(f \in \mathrm{A}(\mathrm{G})\). Let us take \(g = \widetilde{f}\) in formula (21). It is then written as

\[
\int_ {\widehat {\mathrm{G}}} | \mathcal {F} (f) | ^ {2} d \nu = \int_ {\mathrm{G}} | f | ^ {2} d x,\tag{22}
\]

which proves that the measure \(\nu\) is not zero. The measure \(\nu\) is therefore a Haar measure on \(\widehat{G}\).

DEFINITION 4. --- The Haar measure \(\nu\) on \(\widehat{G}\) from Proposition 9 is called the dual measure of the given Haar measure dx on G.

We will often denote by \(d\chi\) or \(d\widehat{x}\) the Haar measure on \(\widehat{G}\) which is dual to the Haar measure dx.

Note. --- Let \(a\) be a real number \(>0\). If one replaces \(dx\) by the measure \(a \cdot dx\), the convolution product of the functions \(f\) and \(g \in \mathrm{L}^{1}(\mathrm{G})\) is replaced by \(a(f*g)\). We have seen (II, p.~209, remark) that \(\mathcal{F}(f)\) is replaced by \(a\mathcal{F}(f)\). Therefore \(\mu_{f}\) is unchanged and \(\nu\) is replaced by \(a^{-1} \cdot \nu\). In particular, the measure \(dx \otimes d\hat{x}\) on \(\mathrm{G} \times \hat{\mathrm{G}}\) is independent of the choice of Haar measure on \(\mathrm{G}\).

Lemma 5. --- The space A(\(\hat{G}\)) is dense in L²(\(\hat{G}\)).

Let \(h\) be an element of \(\mathrm{L}^2(\widehat{\mathrm{G}})\) orthogonal to \(\widehat{\mathrm{A}}(\mathrm{G})\). For \(f\) and \(g\) in \(\mathrm{A}(\mathrm{G})\), one has \(\mathcal{F}(f) \cdot \mathcal{F}(g) = \mathcal{F}(f*g) \in \widehat{\mathrm{A}}(\mathrm{G})\), hence \(h \cdot \overline{\mathcal{F}(f)}\) is orthogonal to \(\mathcal{F}(g)\). Thus, for every \(f \in \mathrm{A}(\mathrm{G})\), the function \(h \cdot \mathcal{F}(f)\) is orthogonal to \(\widehat{\mathrm{A}}(\mathrm{G})\). But \(h \cdot \mathcal{F}(f) \in \mathrm{L}^1(\widehat{\mathrm{G}})\), and \(\widehat{\mathrm{A}}(\mathrm{G})\) is dense in \(\mathcal{C}_0(\widehat{\mathrm{G}})\), so the measure \(h\mathcal{F}(f) \cdot \nu\) is zero, that is, \(h\mathcal{F}(f)\) is \(\nu\)-locally negligible (INT, V, §5, no. 3, cor. 2 of prop. 3). In particular, \(h\) is \(\nu\)-locally negligible on the set \(\Omega_f\) of characters \(\chi\) such that \(\mathcal{F}(f)(\chi) \neq 0\). By corollary 2, it follows that \(h\) is \(\nu\)-locally negligible, hence zero since \(h\) belongs to \(\mathrm{L}^2(\widehat{\mathrm{G}})\). This concludes the proof.

THEOREM 1 (Plancherel). --- The restriction of the Fourier transform to the subspace A(G) of L²(G) extends uniquely to an isometry Φ of L²(G) onto L²(Ĝ).

Moreover, if \(f \in \mathrm{L}^{1}(\mathrm{G}) \cap \mathrm{L}^{2}(\mathrm{G})\), its Fourier transform belongs to \(\mathrm{L}^{2}(\widehat{\mathrm{G}})\) and coincides in \(\mathrm{L}^{2}(\widehat{\mathrm{G}})\) with \(\Phi(f)\).

By formula (22), the restriction of \(\mathcal{F}\) to A(G) is an isometry of the subspace A(G) of L \(^2\) (G) onto the subspace \(\hat{\mathrm{A}}(\mathrm{G})\) of L \(^2\) (\(\hat{\mathrm{G}}\)). Since A(G) is dense in L \(^2\) (G) (cor. 1 of II, p.~212), the Fourier transform extends uniquely to an isometry \(\Phi\) of L \(^2\) (G) onto a closed subspace of L \(^2\) (\(\hat{\mathrm{G}}\)). But since its image contains \(\hat{\mathrm{A}}(\mathrm{G})\), which is dense in L \(^2\) (\(\hat{\mathrm{G}}\)) (lemma 5), the map \(\Phi\) is surjective.

Now let \(f \in \mathrm{L}^{1}(\mathrm{G}) \cap \mathrm{L}^{2}(\mathrm{G})\); let us show that its Fourier transform belongs to \(\mathrm{L}^{2}(\widehat{\mathrm{G}})\). By prop. 8, (iv) of II, p.~211, and the fact that \(\mathrm{A}(\mathrm{G})\) is an ideal of \(\mathrm{L}^{1}(\mathrm{G})\), there exists a filter base B on \(\mathrm{A}(\mathrm{G})\) which converges to f both in \(\mathrm{L}^{1}(\mathrm{G})\) and in \(\mathrm{L}^{2}(\mathrm{G})\). We then have

\[
\Phi (f) = \lim _ {g, \mathfrak {B}} \Phi (g) = \lim _ {g, \mathfrak {B}} \mathcal {F} (g)
\]

in \(\mathrm{L}^{2}(\widehat{\mathrm{G}})\), and \(\mathcal{F}(f)=\lim_{g,\mathfrak{B}}\mathcal{F}(g)\) in \(\mathcal{C}_{0}(\widehat{\mathrm{G}})\). There therefore exists a sequence \((g_{n})\) in \(\mathrm{A}(\mathrm{G})\) such that \(\mathcal{F}(g_{n})\) converges to \(\Phi(f)\) in \(\mathrm{L}^{2}(\widehat{\mathrm{G}})\) and to \(\mathcal{F}(f)\) in \(\mathcal{C}_{0}(\widehat{\mathrm{G}})\). By INT, IV, §3, no. 4, th. 3 and cor. 1, we have \(\mathcal{F}(f)=\Phi(f)\), and in particular \(\mathcal{F}(f)\in\mathrm{L}^{2}(\widehat{\mathrm{G}})\). This completes the proof of the theorem.

We still denote by \(\mathcal{F}\) the isometry of \(\mathrm{L}^2 (\mathrm{G})\) onto \(\mathrm{L}^2 (\widehat{\mathrm{G}})\) defined in Theorem 1, and we call it the Fourier transformation on \(\mathrm{L}^2 (\mathrm{G})\). Likewise, the Fourier cotransform admits a unique isometric extension to \(\mathrm{L}^2 (\mathrm{G})\), still called the Fourier cotransform and denoted \(\overline{\mathcal{F}}\).

COROLLARY. --- Suppose that G is compact and that dx is the normalized Haar measure on G. Then the family of unitary characters of G is an orthonormal basis of L²(G).

Since G is compact, the characters of G belong to \(L^{2}(G)\). For \(\chi\) and \(\xi\) in \(\widehat{G}\), one has

\[
\int_ {\mathrm{G}} \overline {{\langle \chi , x \rangle}} \langle \xi , x \rangle d x = \mathcal {F} _ {\mathrm{G}} (d x) (\chi \xi^ {- 1}),
\]

so the family of characters of G is orthonormal (prop. 6 of II, p.~210). It is moreover total because the scalar product of \(\chi \in \widehat{\mathbf{G}}\) and of \(f \in \mathrm{L}^2(\mathrm{G})\) is equal to \(\mathcal{F}_{\mathrm{G}}(f)(\chi)\), and hence \(f\) is orthogonal to \(\widehat{\mathbf{G}}\) if and only if \(\mathcal{F}_{\mathrm{G}}(f)\) is zero in \(\mathrm{L}^2(\widehat{\mathbf{G}})\), if and only if \(f\) is zero (th. 1).

Remarks. --- 1) Some of the formulas concerning the Fourier transformation on \(\mathrm{L}^{1}(\mathrm{G})\) extend to the Fourier transformation on \(\mathrm{L}^{2}(\mathrm{G})\). In particular, for \(f \in \mathrm{L}^{2}(\mathrm{G})\) and \(\chi \in \widehat{\mathbf{G}}\), one has

\[
\overline {{\mathcal {F}}} (f) = (\chi \mapsto \mathcal {F} (f) (\chi^ {- 1})) = \overline {{\mathcal {F} (\overline {{f}})}}
\]

\[
\mathcal {F} (\varepsilon_ {x} * f) = \eta (x ^ {- 1}) \mathcal {F} (f), \qquad \overline {{\mathcal {F}}} (\varepsilon_ {x} * f) = \eta (x) \mathcal {F} (f)
\]

\[
\mathcal {F} (\chi f) = \varepsilon_ {\chi} * \mathcal {F} (f), \quad \overline {{\mathcal {F}}} (\chi f) = \varepsilon_ {\chi} * \mathcal {F} (f).
\]

If \(\sigma\) is an automorphism of G and \(\Delta\) the modulus of \(\sigma\) (INT, VII, §1, no. 4, def. 4), then for \(f \in \mathrm{L}^2(\mathrm{G})\), one has

\[
\mathcal {F} (f \circ \sigma) = \Delta^ {- 1} \mathcal {F} (f) \circ \widehat {\sigma} ^ {- 1}
\]

in L \(^{2}\) (G).

\begin{enumerate}
\def\labelenumi{\arabic{enumi})}
\setcounter{enumi}{1}
\tightlist
\item
  The formulas
\end{enumerate}

\[
\| \overline {{\mathcal {F}}} (f) \| ^ {2} = \| f \| ^ {2}\tag{23}
\]

for \(f\in \mathrm{L}^2 (\mathrm{G})\), or

\[
\int_ {\mathrm{G}} f (x) \overline {{g (x)}} d x = \int_ {\widehat {\mathrm{G}}} \mathcal {F} (f) (\chi) \overline {{\mathcal {F} (g) (\chi)}} d \chi\tag{24}
\]

for \(f\) and \(g\) in \(\mathrm{L}^2 (\mathrm{G})\), are called “Plancherel formulas”.

PROPOSITION 10. --- Let \(n \geqslant 0\) be an integer and let \(\mathrm{G}_{1}, \cdots, \mathrm{G}_{n}\) be locally compact abelian groups. Let \(\mu_{j}\), for \(1 \leqslant j \leqslant n\), be Haar measures on \(\mathrm{G}_{j}\). Let G be the product group of the groups \(\mathrm{G}_{j}\) for \(1 \leqslant j \leqslant n\). Let \(\beta\) be the isomorphism from \(\widehat{\mathrm{G}}\) onto \(\prod \widehat{\mathrm{G}}_{j}\) of prop. 2 of II, p.~206. The Haar measure on \(\widehat{\mathrm{G}}\) dual to the product Haar measure \(\mu = \mu_{1} \otimes \cdots \otimes \mu_{n}\) on G is identified with the product of the Haar measures \(\widehat{\mu}_{j}\).

Indeed, for every family \((f_{j})\) of nonzero elements of \(\mathcal{L}^{2}(\mathrm{G}_{j})\), the function f on G defined by \((x_{j})\mapsto\Pi f_{j}(x_{j})\) belongs to \(\mathcal{L}^{2}(\mathrm{G})\), and satisfies

\[
\int_ {\mathrm{G}} | f | ^ {2} d \mu = \prod_ {j} \int_ {\mathrm{G} _ {j}} | f _ {j} | ^ {2} d \mu_ {j} = \prod_ {j} \int_ {\widehat {\mathrm{G}} _ {j}} | \mathcal {F} _ {\mathrm{G} _ {j}} (f _ {j}) | ^ {2} d \widehat {\mu} _ {j},
\]

by the Plancherel formula, which proves that the product Haar measure of the \(\widehat{\mu}_{j}\) is identified with the dual Haar measure of \(\mu\).

